\documentclass[11pt,a4paper]{article}
\usepackage[T1]{fontenc}
\usepackage{lmodern,amsmath,amssymb,amsthm,mathtools,booktabs,microtype}
\usepackage[margin=27mm]{geometry}
\usepackage[colorlinks=true,linkcolor=blue,citecolor=blue,urlcolor=blue]{hyperref}
\newtheorem{theorem}{Theorem}[section]
\newtheorem{lemma}[theorem]{Lemma}
\newtheorem{proposition}[theorem]{Proposition}
\newtheorem{corollary}[theorem]{Corollary}
\newcommand{\Q}{\mathbb Q}\newcommand{\Z}{\mathbb Z}
\DeclareMathOperator{\Tr}{Tr}\DeclareMathOperator{\Nm}{N}
\setlength{\emergencystretch}{2em}
\title{An Explicit Defining Polynomial in Twenty-Seven Variables}
\author{}\date{10 November 2026}
\begin{document}\maketitle
\begin{abstract}
Under the specified normalized elliptic, modular, pole-realization and
prime-pattern inputs, we construct integer polynomials in twenty-seven
total variables whose rational input projection is exactly the integers.
An explicit sufficient coefficient parameter is $L=125001$.
Three comparison scales control the two rational components of a sum
of four elliptic abscissas, so the height argument uses only effective
constants. Scalar and comparison divisibility conditions share three
pole predicates. A clustered choice of the comparison primes permits
one additional congruence on the common elliptic index and makes linear
scalar normalization sufficient. Nine quartic roots and three quadratic
roots are eliminated with one separating coordinate. An exact integer
certificate verifies the finite exceptional set and every numerical
inequality in the sufficient threshold.
\end{abstract}

\section{Statement and arithmetic inputs}

A polynomial $F\in\Z[X,Y_1,\ldots,Y_m]$ defines $\Z$ over $\Q$
when
\[
 x\in\Z\quad\Longleftrightarrow\quad
 \exists\mathbf y\in\Q^m:\ F(x,\mathbf y)=0
 \qquad(x\in\Q).
\]
All variable counts include the input $X$.

\begin{theorem}\label{thm:main}
Under the arithmetic inputs below there is a computable family
\[
 F_L\in\Z[X,Y_1,\ldots,Y_{26}],\qquad L\ge3\text{ odd},
\]
given by a finite integer coefficient specification. Every odd
$L\ge125001$ defines $\Z$ over $\Q$. In particular,
$F_{125001}$ is a fully specified defining polynomial.
The integer $L$ is a fixed coefficient parameter and is not a
quantified variable. The total degree is at most
$99950805468739574201057280<2^{87}$, uniformly in $L$.
For $L=125001$, the maximum absolute integer coefficient is less
than $2^{2^{102}}$. No additional $abc$ or Szpiro hypothesis is used.
\end{theorem}

Write $h$ for the absolute logarithmic Weil height, and put
\[
 K=\Q(\sqrt2),\quad d=3-\sqrt2,\quad D=d^2=11-6\sqrt2,
 \quad E:y^2=x^3-Dx,
\]
\[
 P=\left(\frac{57-28\sqrt2}{8},
          \frac{-660+417\sqrt2}{32}\right).
\]
The nonidentity automorphism of $K$ is denoted by $\sigma$.
The following four arithmetic inputs have the scope established in
\cite{Dependencies,Source}.
\begin{enumerate}
\item $2E(K)=\Z P$, the point $P$ is non-torsion, and
$h(x([n]P))\ge n^2-6$ for every nonzero integer $n$.
\item The full normalized compatible modular realization for the
five contacts is the one in~\cite{Multiplicity}. Its proved consequence is
\begin{equation}\label{eq:baseheight}
 h(q)\le H M_5(q)^{7/3},\qquad H=2^{221234},
\end{equation}
for every rational $q$ outside $0,1,2401,-343/32$. If $q=u/v$
is reduced with $v>0$, $M_5(q)$ is the product of primes
$p>76176$ having positive odd valuation in at least one of
\[
                 v,\quad u,\quad u-v,\quad u-2401v,
                         \quad32u+343v.
\]
The modular realization, including the normalized compatible
representation used in its multiplicity theorem, is retained as
an explicit arithmetic input.
\item For $b\ne0$ and $B_0=17b$, the pole predicate is
\begin{equation}\label{eq:pole}
 \begin{split}
 p_0^2&=B_0(1-e)(1+3e),\\
 B_0q_0^2&=(s^2-B_0e(B_0-b^2))(s^2+3B_0e(B_0-b^2)).
 \end{split}
\end{equation}
A rational solution excludes negative odd valuations of $b$ outside
$\{2,3,17\}$. A solution with $p_0q_0s\ne0$ exists whenever
\begin{equation}\label{eq:prepared}
 \begin{gathered}
 0<b<17,\qquad b\text{ has square reduced denominator},\\
 b\in\Q_2^{\times2}\cap\Q_3^{\times2}\cap\Q_{17}^{\times2},
 \qquad v_2(b)\ge3,\quad v_3(b)\ge1,\quad v_{17}(b)\ge2.
 \end{gathered}
\end{equation}
The restricted-converse and descent inputs supporting this assertion
are those identified in~\cite{Dependencies,PointSource}.
\item The simultaneous prime-pattern input applies to any fixed
finite list of rational contact coefficients. It supplies positive
prime numerators, the prescribed Gaussian generators and conjugate
reduction orders, arbitrary fixed denominator divisibility, arbitrary
fixed avoidance, and arbitrarily large real slope. Its precise
consequence of Sun's theorem is stated in~\cite{Dependencies}.
\end{enumerate}

The effective abscissa-sum estimate in~\cite{Sums} is proved from the
first input. We will use it explicitly in Section~\ref{sec:heights}.
No ineffective aggregate trace-height bound enters the argument.

\section{The rational source}

Every displayed rational expression is restricted to its domain.
All denominator guards are retained when clearing fractions. We also
require the comparison slope $\lambda$ and every pole input to be
nonzero.

\subsection{Nine elliptic points}

Retain the nine points
\[
 A,S,T_1,Z,H_0,U_A,U_S,U_T,U_Z\in2E(K)\setminus\{O\}.
\]
For each point use two rational coordinates $h_0,h_1$ and the finite
chart
\[
 h=h_0+h_1\sqrt2,\qquad x=h^2,\qquad j^2=h^4-D,\qquad y=hj.
\]
Under the group input these are precisely the nonzero multiples of
$P$. Indeed the two-isogeny descent map shows that the square-abscissa
subgroup is $2E(K)$; the squareclass and sign checks for the displayed
curve are proved in~\cite{SquareCharts}.

For clarity, the one-root quartic chart for $j$ is as follows. Write
\[
 A_0=h_0^4+12h_0^2h_1^2+4h_1^4-11,\qquad
 B_0=4h_0h_1(h_0^2+2h_1^2)+6,
\]
\[
 c_0=2A_0+B_0,\quad d_0=-A_0-B_0,
\]
\[
 A'=6A_0^3-9A_0B_0^2-4B_0^3,\qquad
 B'=-4A_0^3+6A_0B_0^2+3B_0^3.
\]
Adjoin one root $Y$ of
\begin{equation}\label{eq:pointquartic}
 Y^4-4A'Y^2+8(B')^2=0,\qquad
 j=\frac{Y/2+(B'/Y)\sqrt2}{c_0+d_0\sqrt2}.
\end{equation}
This chart is complete and separable at every rational retained-base
specialization~\cite{QuadraticCharts,SquarePredicate}. Its discriminant is
\[
 2^{15}B'^2(2A_0^2-B_0^2)^4(A_0^2-2B_0^2)^2\ne0.
\]
The nonvanishing follows from the norm obstruction to $h^4=D$,
the identity
$A'^2-2B'^2=(2A_0^2-B_0^2)^2(A_0^2-2B_0^2)$,
and the Eisenstein cubic $4T^3-6T-3$ excluding $B'=0$.

Write the integer labels of $(A,S,T_1,Z)$ as $(n,a,t,b)$, and set
\[
 c=|a-1|,\qquad B=\max(|n|,|a|,|t|,|b|).
\]
The affine formula for $S-P$ includes its equal-abscissa guard, so
$a\ne\pm1$ and $c>0$. Put
\[
 Q_0=P,\quad (Q_1,Q_2,Q_3,Q_4)=(A,S,T_1,Z),\quad U_0=H_0,
\]
\[
 (U_1,U_2,U_3,U_4)=(U_A,U_S,U_T,U_Z),\qquad V_i=U_i-Q_i.
\]
All differences use the usual guarded affine addition formulas.

\subsection{Comparison and scalar functions}

Write $f(Q)=1/x(Q)$ and $z(Q)=-x(Q)/y(Q)$. The order-four
formal-logarithm truncation is $z(Q)$ itself. Set
\[
 L_i=\frac{z(U_i)}{z(H_0)},\qquad
 L_i'=\sigma\left(\frac{z(V_i)}{z(V_0)}\right)
 \quad(1\le i\le4),
\]
and
\[
 \mathcal F(u,v,w)=v^2-u(u-2)vw+(u-2)^2w^2-1,
 \qquad \beta=216X^4+450459.
\]
The sixteen comparison functions, in order, are
\begin{equation}\label{eq:comparisons}
 \begin{gathered}
 (f(U_i))_{i=0}^4,\qquad (\sigma f(V_i))_{i=0}^4,\\
 (L_i-L_i')_{i=1}^4,\qquad L_1-\beta,\qquad
                              \mathcal F(L_2,L_3,L_4).
 \end{gathered}
\end{equation}
Denote them by $C_1,\ldots,C_{16}$, and put
\[
 m_i=1\quad(1\le i\le10),\qquad
 m_i=3\quad(11\le i\le16).
\]
Let $\eta_1,\ldots,\eta_6\in\Q$ be the two rational coefficients
in the basis $1,\sqrt2$ of each of
\begin{equation}\label{eq:scalar-functions}
             f(H_0),\qquad f(S-P),\qquad z(S-P)/z(H_0),
\end{equation}
in that order.

For fixed odd $L\ge3$, define
\begin{equation}\label{eq:rho}
 \begin{gathered}
 \ell=76249,\qquad
 r=\beta(\beta-1)(\beta-2401)(32\beta+343)(2\beta-1),\\
 C_L=17(\ell^{(L+1)/2}+1),\qquad
 \rho=\frac{102^2\ell^Lr}{(1+C_Lr)^2}.
 \end{gathered}
\end{equation}
The scalar normalization theorem~\cite{Normalizer} gives
\begin{equation}\label{eq:forcedrho}
 v_\ell(\rho)=
 \begin{cases}L,&v_\ell(X)\ge0,\\L-20v_\ell(X),&v_\ell(X)<0.
 \end{cases}
\end{equation}
This valuation is positive and odd. The same theorem gives, for real $X$,
\begin{equation}\label{eq:smallrho}
                  0<\rho<\frac{36}{\ell r}<1.
\end{equation}
For integral $X$, the reduced denominator of $\rho$ is a square;
its squareclasses and depths at $2,3,17$ satisfy~\eqref{eq:prepared},
and $v_2(\rho)=4$.

\subsection{Three scales and three pole inputs}

Put
\[
 \Theta=\{0,-1/2,-1,1/2399,-32/407\},
\]
\[
 \Theta_j=(4j+2+\Theta)\cup(4j+4+\Theta),\qquad
 \gamma_j=-j-1\quad(j=0,1,2).
\]
All thirty contacts are positive. The three contact-and-pole sets
are pairwise disjoint. Write
\[
 \mathcal S=\sum_{i=1}^4x(Q_i)=S_0+S_1\sqrt2,
 \qquad U=S_0,\qquad V=S_0+S_1,
\]
and define
\begin{equation}\label{eq:scales}
 A_0=1,\qquad A_1=(1+\rho U^2)^2,\qquad
 A_2=(1+\rho V^2)^2,\qquad t_j=A_j\lambda.
\end{equation}
These scales are positive rational squares. For $j=0,1,2$, set
\begin{equation}\label{eq:contactfunctions}
 \delta_j=t_j-\gamma_j,\qquad
 R_j=\frac{\prod_{\theta\in\Theta_j}(t_j-\theta)}{\delta_j^{10}},
 \qquad w_j=\rho R_j.
\end{equation}
Require every contact and $\delta_j$ to be nonzero.

For nonzero rational $q$ put $\Phi_q(v)=qv/(1+qv)$ on its
nonzero-denominator domain. Define
\begin{equation}\label{eq:weights}
 G_{j,k}=R_j^3\Phi_{R_j}(\eta_k),
\end{equation}
\begin{equation}\label{eq:W}
 W_j=\prod_{i=1}^{16}
       \left(1+w_j\left(\frac{C_i}{\delta_jw_j^{m_i}}\right)^2\right)
     \prod_{k=1}^{6}
       \left(1+w_j\left(\frac{G_{j,k}}{w_j^3}\right)^2\right),
\end{equation}
\begin{equation}\label{eq:poleinputs}
                      b_j=\frac{4\rho R_j}{\Tr_{K/\Q}(W_j^{-1})^2}.
\end{equation}
Each $b_j$ receives the pole predicate~\eqref{eq:pole}. Its first
conic uses two retained rational parameters $\tau_j,s_j$ and
\[
 e_j=\frac{\tau_j^2-B_j}{\tau_j^2+3B_j},\qquad
 p_j=\frac{-4B_j\tau_j}{\tau_j^2+3B_j},\qquad B_j=17b_j.
\]
The second equation is a single nonzero quadratic root after rational
square clearing. The nonzero-witness form of the inherited converse
makes this conic chart complete~\cite{PointSource}.

Let $\Sigma$ be the primes at most $76176$, together with every
prime dividing a numerator or denominator of a contact, a pole, or
a nonzero difference between two such values. Exact factorization gives
\begin{equation}\label{eq:exceptional}
 \Sigma\setminus\{p:p\le76176\}
 =\{149027,272011,1875611,3828397,3982747\}.
\end{equation}
Their product, denoted by $E_*$, satisfies $E_*<2^{100}$, and
$\ell\notin\Sigma$. The checker in Section~\ref{sec:certificate}
verifies these assertions independently by integer trial division.

\section{Local divisibility}

At primes outside $\Sigma$ all valuations on completions of $K$
are integer valued and normalized by $v(p)=1$. A zero test function
has valuation $+\infty$.

\begin{lemma}\label{lem:trace}
Suppose $\alpha=v_p(w_j)$ is odd and $b_j$ has a pole certificate.
Then $\alpha>0$. At some place of $K$ above $p$, every factor
of $W_j$ has valuation greater than $-\alpha/2$.
\end{lemma}
\begin{proof}
The nonzero second summand of a factor $1+w_jq^2$ has odd
valuation, so that factor has valuation at most zero. A zero second
summand gives the unit factor $1$. Hence $W_j^{-1}$ is integral at
every place above $p$, and its trace is integral. The pole obstruction
and~\eqref{eq:poleinputs} give
\[
 0\le\alpha-2v_p(\Tr(W_j^{-1})),\qquad
 v_p(\Tr(W_j^{-1}))\le(\alpha-1)/2.
\]
Thus $\alpha>0$. At least one local value of $W_j^{-1}$ has
valuation at most $(\alpha-1)/2$; otherwise so would its trace.
At that place the sum of the nonpositive valuations of the factors
of $W_j$ is greater than $-\alpha/2$. Each summand is therefore
greater than $-\alpha/2$.
\end{proof}

\begin{lemma}\label{lem:unit}
If $v_p(\rho)$ is odd, then at least one of $R_0,R_1$ is a unit.
\end{lemma}
\begin{proof}
A balanced function $R_j$ is a unit whenever $v_p(t_j)\ne0$,
because all contacts and poles are units. If $U=0$, then
$A_1=1$, and disjointness of the two reduced contact-and-pole sets
proves the assertion directly. Otherwise write
$\kappa=v_p(\rho U^2)$, which is odd. If $\kappa>0$, then
$A_1\equiv1\pmod p$. For nonunit $\lambda$, both arguments
are nonunits and both contact functions are units. For unit
$\lambda$, the arguments have the same residue; the disjointness
of the two contact-and-pole sets modulo $p$ prevents both functions
from being nonunits. If $\kappa<0$, then $v_p(A_1)=2\kappa\ne0$.
The arguments $t_0$ and $t_1$ cannot both be units, so at least one
of their balanced functions is a unit.
\end{proof}

\begin{lemma}\label{lem:scalar}
If $\mu=v_p(\rho)$ is odd, then $\mu>0$ and
\[
                           v_p(c)\ge2\mu+1\ge3.
\]
In particular
\begin{equation}\label{eq:cforced}
               c\ge\ell^{2L+1},\qquad 4L\log\ell\le2\log c.
\end{equation}
\end{lemma}
\begin{proof}
Choose a unit $R_j$ by Lemma~\ref{lem:unit}. Lemma~\ref{lem:trace}
forces $\mu>0$. Applied to each scalar factor, it gives
\[
                           v(G_{j,k})>9\mu/4.
\]
For a unit $R_j$, positive valuation of $\Phi_{R_j}(\eta_k)$
implies positive valuation of $\eta_k$, and the two valuations
are equal. Therefore every rational scalar component has valuation
$>9\mu/4$, hence at least $2\mu+1$.

Both $H_0$ and $S-P$ are then in the formal group at every place
above $p$. If $H_0=[h]P$, the full formal-logarithm ratio is
$(a-1)/h$. The quotient of a formal parameter and its logarithm
is a unit. The third pair of scalar components consequently gives
\[
 v_p(a-1)-v_p(h)
       =v(z(S-P)/z(H_0))\ge2\mu+1.
\]
Since $h$ is an integer, this proves the asserted bound for $c$.
Apply it at $\ell$, using~\eqref{eq:forcedrho}, to obtain
\eqref{eq:cforced}.
\end{proof}

\begin{lemma}\label{lem:comparison}
Let $v_p(R_j)=\nu>0$ be odd and suppose $v_p(\rho)$ is even.
At some place above $p$, the first ten functions in
\eqref{eq:comparisons} have positive valuation and the last six
have valuation at least three.
\end{lemma}
\begin{proof}
At such a contact prime $\delta_j$ is a unit. The integer
$\alpha=v_p(\rho)+\nu$ is odd. By Lemma~\ref{lem:trace},
$\alpha>0$ and the comparison factors at one place satisfy
\[
 \alpha+2\bigl(v(C_i)-m_i\alpha\bigr)>-\alpha/2
\]
whenever the second summand is not integral. The same resulting
lower bound holds if it is integral. Thus
\[
                   v(C_i)>(m_i-3/4)\alpha\ge m_i-3/4,
\]
which gives $v(C_i)\ge m_i$ by integer rounding.
\end{proof}

For a point in the first formal neighborhood at $p>76176$, one has
\begin{equation}\label{eq:formal}
 v(\log_E(Q)-z(Q))\ge5v(z(Q)),\qquad
                         z(Q)/\log_E(Q)\in1+p^4\mathcal O.
\end{equation}
This follows from the formal logarithm of $y^2=x^3-Dx$, whose
first omitted term has parameter order five; the integral-denominator
estimate at these primes is proved in~\cite{PointSource,RawSource}.
Positive valuation of $1/x(Q)$ is equivalent to membership in the
first formal neighborhood.

Suppose Lemma~\ref{lem:comparison} applies, and write
$U_i=[u_i]P$, $H_0=[h]P$, and
$(n_1,n_2,n_3,n_4)=(n,a,t,b)$. The exact logarithm ratios are
\[
             r_i=u_i/h,\qquad r_i'=(u_i-n_i)/(h-1).
\]
At least one of $h,h-1$ is a $p$-unit, so its ratio is integral.
The cross comparison and~\eqref{eq:formal} make the other truncated
ratio integral as well. Both relative errors are then absolute
errors, and modulo $p^3$ one has $L_i\equiv r_i\equiv r_i'$.
Multiplying by $h,h-1$ and subtracting gives
\[
                             L_i\equiv n_i\pmod{p^3}.
\]
Define the integer residual $\Delta=\mathcal F(a,t,b)$. If
$\beta=u/v$ is reduced with $v>0$, the direct target also makes
$\beta$ integral at $p$. Consequently
\[
                         p^3\mid\Delta,\qquad p^3\mid u-nv.
\]

Let $M_j$ be the product of primes outside $\Sigma$ at which $R_j$
has positive odd valuation. At a prime in $M_j$, an odd scalar
valuation invokes Lemma~\ref{lem:scalar}; an even scalar valuation
invokes Lemma~\ref{lem:comparison}. Therefore, for either nonzero
integer residual $N=\Delta$ or $N=u-nv$,
\begin{equation}\label{eq:divisibility}
                         M_j^3\mid Nc\quad(j=0,1,2).
\end{equation}
No prime shared by the scalar and comparison coefficients is omitted.

\section{Effective height bounds}\label{sec:heights}

Set
\begin{equation}\label{eq:Hbar}
                         \overline H=2^{221470}.
\end{equation}
All constants in this section are independent of $X$ and $L$.

\subsection{Scalar contacts and a coprime divisibility}

Let $M_\beta^\circ$ be the product of primes in $M_5(\beta)$
outside $\Sigma$. Such a prime has positive odd valuation in $\rho$.
Indeed, a negative valuation of $\beta$ is a multiple of four,
so its denominator cannot contribute an odd contact. At a finite
contact exactly one of the five factors of $r$ has positive odd
valuation, and $1+C_Lr$ is a unit. The additional factor $2\beta-1$
is a unit at these four finite contacts outside $\Sigma$.
Lemma~\ref{lem:scalar} gives
\[
                         (M_\beta^\circ)^3\mid c.
\]
The prime $\ell$ does not contribute to $M_5(\beta)$. For
$\ell$-integral $X$, enumeration of fourth-power residues excludes
$\beta\equiv0,1,2401,-343/32\pmod\ell$; for nonintegral $X$,
$v_\ell(\beta)$ is a multiple of four. The independent certificate
in Section~\ref{sec:certificate} checks the finite enumeration.
Combining this coprimality with~\eqref{eq:cforced} yields the stronger
product divisibility
\begin{equation}\label{eq:coprime}
               \ell^{2L+1}(M_\beta^\circ)^3\mid c.
\end{equation}

Removing the five additional exceptional primes from a contact
radical changes it by a factor at most $E_*$. Therefore
\eqref{eq:baseheight} and $E_*<2^{100}$ give
\begin{equation}\label{eq:scalarheight}
 h(\beta)\le\overline H(M_\beta^\circ)^{7/3}
       \le\overline H\left(\frac{c}{\ell^{2L+1}}\right)^{7/9}
       \le\frac{\overline H}{\ell^{2L+1}}c.
\end{equation}
The last inequality uses $c/\ell^{2L+1}\ge1$. The factor
$\ell^{2L+1}$ is retained before applying the height bound; it is
not absorbed into an unspecified constant.

The scalar coefficient also has the elementary estimate
\begin{equation}\label{eq:rhoheight}
                    h(\rho)\le15h(\beta)+2L\log\ell+128.
\end{equation}
To prove it, put $K_0=2\cdot2402\cdot375\cdot3=5404500$.
The elementary height inequalities give
\[
 h(r)\le5h(\beta)+\log K_0,
\]
\[
 h(\rho)\le\log102^2+L\log\ell+3h(r)+2\log C_L+2\log2.
\]
Since $C_L\le34\ell^{(L+1)/2}$, the remaining constant is
\[
 3\log K_0+2\log34+\log\ell+\log102^2+2\log2<114<128.
\]
The displayed bound $114$ follows already from the binary lengths
of the five positive integers and $\log2<1$.

\subsection{Ten-contact heights}

For $s\in\{2,4,6,8,10,12\}$ the rational map
\[
                         q_s=2+\frac1{t-s}
\]
sends $s+\Theta$ to the five contacts
$\infty,0,1,2401,-343/32$. Its inverse has the integral matrix
\[
                         \begin{pmatrix}s&1-2s\\1&-2\end{pmatrix},
\]
whose determinant is $-1$ and whose largest absolute row sum is
at most $35$. Thus $h(t)\le h(q_s)+\log35$.
Outside $\Sigma$, this map preserves contact multiplicities, and
the odd radical of $q_s$ is exactly the odd radical of the five
translated contacts. Denote that radical by $M_s$.

For each $j$, the two five-contact sets are disjoint modulo every
prime outside $\Sigma$. The denominator $\delta_j^{10}$ contributes
only even valuations. Hence
\[
                         M_{4j+2}M_{4j+4}=M_j.
\]
If $h(t_j)>\log35$, applying~\eqref{eq:baseheight} to both
five-contact sets and multiplying gives
\[
 \bigl(h(t_j)-\log35\bigr)^2
                  \le H^2E_*^{14/3}M_j^{7/3}.
\]
The same final bound is immediate if $h(t_j)\le\log35$.
Since
\[
 \log35+H E_*^{7/3}<2^{221470}=\overline H,
\]
we conclude that
\begin{equation}\label{eq:tenheight}
                         h(t_j)\le\overline H M_j^{7/6}
                         \qquad(j=0,1,2).
\end{equation}
All excluded primes and the height cost of translation have now
been included explicitly.

\subsection{Two rational components of the abscissa sum}

The effective sum-height theorem~\cite{Sums} gives
\[
                         B^2\le2h(\mathcal S)+80.
\]
The variables $U=S_0$ and $V=S_0+S_1$ satisfy
\[
                         \mathcal S=(1-\sqrt2)U+\sqrt2 V.
\]
At finite places both coefficients are integral. The product of
the two real coefficient sums is
$(2\sqrt2-1)(2\sqrt2+1)=7$. It follows that
\[
 h(\mathcal S)\le h(U)+h(V)+\tfrac12\log7,
 \qquad B^2\le2h(U)+2h(V)+82.
\]
This uses $\log7<2$.

For $W=U,V$, respectively, the two nontrivial scales satisfy
\[
 h((1+\rho W^2)^2)
       \ge4h(W)-2h(\rho)-2\log2.
\]
The nonzero slope gives $A_j=t_j/t_0$. Combining these inequalities
with $2\log2<2$ yields
\begin{equation}\label{eq:sumheight}
 B^2\le h(t_0)+\tfrac12h(t_1)+\tfrac12h(t_2)
                       +2h(\rho)+84.
\end{equation}
Consequently, for either nonzero residual $N=\Delta$ or $N=u-nv$,
\eqref{eq:divisibility}, \eqref{eq:rhoheight}, and
\eqref{eq:tenheight} give the effective common inequality
\begin{equation}\label{eq:commonheight}
 B^2\le2\overline H|Nc|^{7/18}
                      +30h(\beta)+4L\log\ell+340.
\end{equation}

\section{An explicit sufficient coefficient parameter}

Assume throughout this section that $L\ge125001$ is odd.
As $\ell>2^{16}$, equation~\eqref{eq:cforced} gives
\begin{equation}\label{eq:binaryforced}
 c\ge\ell^{2L+1}>2^{4000048}.
\end{equation}
In particular $\ell^{2L+1}>120\overline H$, so
\eqref{eq:scalarheight} implies
\begin{equation}\label{eq:scalarlinear}
                         h(\beta)\le c/120.
\end{equation}
Using $4L\log\ell\le2\log c$, the last three terms of
\eqref{eq:commonheight} are bounded above by
\[
                       c/4+2\log c+340.
\]
For $c\ge100$, this is less than $(c-1)^2/2\le B^2/2$.
For example, use $\log c\le c$: the polynomial
$(c-1)^2/2-9c/4-340$ is positive at $100$ and increasing thereafter.
Thus these three terms can be absorbed into $B^2/2$.

\subsection{Vanishing of the Lucas residual}

If $\Delta\ne0$, the elementary bounds $B>4$, $c\le2B$ and
$|\Delta|\le5B^4$ give $|\Delta c|\le10B^5$.
Equation~\eqref{eq:commonheight} then implies
\[
 B^{1/18}\le4\overline H10^{7/18},
 \qquad
 B\le4^{18}10^7\overline H^{18}<2^{3986520}.
\]
Here $10^7<2^{24}$ gives the last strict inequality. This contradicts
$B\ge c-1$ and~\eqref{eq:binaryforced}. Therefore
\begin{equation}\label{eq:Delta0}
                              \mathcal F(a,t,b)=0.
\end{equation}

We recall the elementary growth consequence of this equation,
including its proof. If $|a|>7$ and $b\ne0$, then
\begin{equation}\label{eq:Lucasgrowth}
                              \log B>c-4.
\end{equation}
Indeed, put $A=|a|$ and $C=|(a-2)b|$. Equation~\eqref{eq:Delta0}
implies that $T=|t|$ and $C$ are positive integer solutions of
$T^2-ATC+C^2=1$. Their signs before taking absolute values must
agree after multiplying $(a-2)b$ by $\operatorname{sgn}(a)$.
For the recurrence $U_0=0$, $U_1=1$,
$U_{j+1}=AU_j-U_{j-1}$, every positive solution is an unordered
pair of adjacent terms. To see this, order its coordinates
$M\ge m>0$ and replace them by
\[
                         \left(m,Am-M\right),\qquad
               Am-M=(m^2-1)/M\in[0,m).
\]
This preserves the equation and descends to $(1,0)$.
Hence $C=U_j$ for some $j\ge1$ and $T\ge U_{j-1}$.
Writing $d=|a-2|$, the recurrence modulo $d$ gives
$U_j\equiv j$ if $a>0$ and
$U_j\equiv(-1)^{j-1}j$ if $a<0$. Since $d\mid C$,
we have $j\ge d\ge A-2$. Also
$U_r>(A-1)^{r-1}$ for $r\ge2$, by induction.
It follows that
\[
 \log T>(j-2)\log(A-1)\ge(A-4)\log7>A-3\ge c-4,
\]
proving~\eqref{eq:Lucasgrowth}. Its hypotheses hold here because
$c$ satisfies~\eqref{eq:binaryforced} and all root labels are nonzero.

\subsection{Integrality of the input}

Suppose $\beta\ne n$. For $\beta=u/v$ reduced with $v>0$, put
$N=u-nv\ne0$. Then
\[
                         |Nc|\le4e^{h(\beta)}B^2.
\]
After the same absorption, equation~\eqref{eq:commonheight} gives
\[
 B^{11/9}\le4\overline H4^{7/18}e^{(7/18)h(\beta)}.
\]
Taking logarithms and using $\log2<1$ yields
\begin{equation}\label{eq:inputlog}
 \log B<181205+\frac7{22}h(\beta)
          \le181205+\frac7{2640}c<c-4.
\end{equation}
The last inequality holds already for $c\ge10^6$; its two sides
have constant slopes and it holds strictly at $10^6$.
This contradicts~\eqref{eq:Lucasgrowth}, so $\beta=n\in\Z$.
If $X=u_0/v_0$ is reduced with $v_0>0$, integrality of
$216X^4+450459$ implies $v_0^4\mid216=2^3 3^3$.
Thus $v_0=1$ and $X\in\Z$. This proves soundness for every
fixed odd $L\ge125001$.

\section{Completeness}

Fix an integer input $X$ and any fixed odd $L\ge3$.
The construction below gives a rational source witness. Its order
of choices is relevant: the prime-pattern input is applied only
after the root labels, scales, neighborhoods and finite avoidance
conditions have been fixed.

\subsection{Scalar preparation and fixed neighborhoods}

Let $\mathcal P$ consist of the finite prime support of $\rho$,
$\Sigma$, and $2,3,17$. Choose $H_0=[h_0]P$, with positive
$h_0>2$ divisible by $64$, in sufficiently deep formal subgroups
at the finite preparation places and sufficiently close to the
identity at both real embeddings. Such real returns exist in the
compact closure of the cyclic subgroup after any finite index
restriction. Now choose a positive $c$ divisible by $6$ and any
additional required finite modulus, with $[c]P$ still closer to
the identity at those places. Put $a=c+1$ and $S=[a]P$.
The six scalar components in~\eqref{eq:scalar-functions} can
thereby be made arbitrarily small simultaneously. For the ratio,
$H_0$ has already been fixed with nonzero local parameter, and
$[c]P$ can approach the identity as closely as needed. Deep formal
subgroups and continuity give this assertion also at ramified
and bad primes, as in the scalar realization of~\cite{Source}.

With $\rho$ fixed, make the scalar factors in~\eqref{eq:W}
arbitrarily close to $1$ whenever $R_j$ is close to $1$ at
a finite preparation prime, and whenever $1/2<R_j<2$ over the
reals. In the latter case also require $|\eta_k|<1/4$, so
$1+R_j\eta_k>1/2$ for all six denominators.

Set $n=\beta$. Define
\[
 U_0=0,\quad U_1=1,\quad U_{m+1}=aU_m-U_{m-1},\qquad
 k=2(a-2),\quad b=U_k/(a-2),\quad t=U_{k+1}.
\]
The recurrence modulo $a-2$ gives $U_k\equiv k\equiv0$, so
$b$ is integral. Its usual determinant identity gives
$\mathcal F(a,t,b)=0$. Since $a$ is odd and $a\equiv1\pmod3$,
the sequence modulo two has period $0,1,1$ and $k\equiv1\pmod3$.
Thus $n,a,t,b$ are positive and odd. The root points and all three
positive square scales are now fixed.

Temporarily use $H=[h]P$, $U_i=[hn_i]P$. Restrict $h$ to a
fixed congruence class at the finite preparation places and to a
compact real neighborhood preserving all scalar bounds. Shrink
these neighborhoods so that all raw comparison functions remain
bounded and all their original charts are defined. The relevant
multiplication and translation maps have only finitely many poles;
at an integer $h>2$ their equal-abscissa failures are absent because
$n_i\ne0$. There are no comparison normalization denominators.
Take the initial modulus to be divisible by $64$ and by the order
of the component group of the compact real closure of $\Z P$.
That component group is finite, since the closure is a compact
real Lie subgroup.

\subsection{A clustered prime choice and the common index}

Apply the rational-contact version of the prime-pattern input to
$-\theta/A_j$, for $j=0,1,2$ and $\theta\in\Theta_j$.
Repeated coefficients are listed once, with all incidences retained.
It provides
\[
                         \lambda=u/v,\qquad v=D_0^2w^2,
\]
with every
$p_i=u-(\theta/A_j)v$ a prescribed positive prime. Require these
primes to avoid every fixed preparation prime, every numerator or
denominator of a scale, every relevant contact difference, and
all primes where a coordinate or the nonzero local parameter of
$P$ or $S-P$ is not a unit. These are finite conditions, fixed
before prime selection. Prescribe enough divisibility of $w$ to
make $\lambda$ arbitrarily large in absolute value at all finite
preparation places, and choose it arbitrarily large and positive
at the real place. The rational-contact and avoidance consequences
are those established in~\cite{Dependencies}.

As $\lambda$ tends to infinity at any preparation place, $R_j$
tends to $1$ and $\delta_j^{-1}$ to zero. The scalar factors
are close to $1$ by preparation. The comparison factors tend
uniformly to $1$ on the bounded comparison neighborhoods. Thus
all required finite and real closeness conditions can be imposed
by fixed requirements before invoking the prime-pattern input.

Impose additionally
\begin{equation}\label{eq:cluster}
 \lambda>4\max_{j,\theta\in\Theta_j}(\theta/A_j),\qquad u\ge25.
\end{equation}
All contacts are positive, so every selected prime lies in
$(3u/4,u)$. The full order of either conjugate reduction of $E$
at any selected prime $p_k$ is even and bounded by
\[
 p_k+1+2\sqrt{p_k}<u+1+2\sqrt u<3u/2.
\]
The evenness follows from the rational two-torsion. For $u\ge25$,
$\sqrt u\le u/5$ and $1<u/10$, proving the last inequality.
If a selected odd prime $p_i$ divided one of these even group
orders, that order would be at least $2p_i>3u/2$, a contradiction.
Consequently every selected prime is coprime to every prescribed
annihilating reduction modulus.

The inherited conjugate-order construction imposes the primary
annihilation on $h$, the opposite annihilation on $h-1$, and the
initial fixed congruence on $h$. Its odd coprimality and two-adic
compatibility are proved in~\cite{Dependencies}. The preceding
paragraph permits the additional simultaneous congruences
\begin{equation}\label{eq:hcongruence}
                              h\equiv0\pmod{\prod_i p_i}.
\end{equation}
All selected primes also avoid the initial fixed modulus.

This refined arithmetic progression still meets the prepared real
neighborhood. Indeed, let $G$ be the compact closure of $\Z P$
and $G^\circ$ its identity component. If the initial modulus $m$
annihilates $G/G^\circ$, then $mG=G^\circ$: containment is clear,
and multiplication by $m$ on the compact torus $G^\circ$ is
surjective. For every positive integer $r$ one likewise has
$mrG=G^\circ$. Hence every compatible refinement of the index
congruence has the same closure in the prepared component.
Choose an index $h>2$ in the refined progression and real
neighborhood. Set $H_0=[h]P$ and $U_i=[hn_i]P$.

\subsection{The linear scalar tests at the selected primes}

At a selected prime $p$ for detector $j$, one has
$v_p(R_j)=1$, while $\rho$ and $\delta_j$ are units. The
reduction-order conditions put the required $U_i$ and $V_i$
in their first formal neighborhoods. The first ten comparison
functions therefore have valuation at least two. Their exact
logarithm ratios are the integers $n_i$; equation~\eqref{eq:formal}
gives valuation at least four for the last six comparisons.
At the designated place every comparison second summand in
\eqref{eq:W} thus has valuation at least three.

We verify separately the six linear scalar denominators. At the
primary place, $[h/p]P$ is in the formal group: division by $p$
preserves reduction annihilation because the reduction modulus
is coprime to $p$. Since $H_0=[p]([h/p]P)$, its local parameter
has valuation at least two. At the opposite place, $H_0$ reduces
to $P$, whose coordinates and parameter are units by avoidance.
The selected primes split in $K$, as provided by the prime pattern.
The two rational components of $f(H_0)$ are units: its two local
values have respectively positive and zero valuation. The two
components of $f(S-P)$ are integral. Finally, the local values
of $z(S-P)/z(H_0)$ have valuations at most $-2$ and exactly zero;
its two rational components consequently both have valuation at
most $-2$. Here $2$ and $\sqrt2$ are units, so the sum and
scaled difference of two values with unequal valuations preserve
the smaller valuation.

Every rational component $\eta_k$ thus has either nonnegative
valuation or valuation at most $-2$. Since $v_p(R_j)=1$,
$v_p(R_j\eta_k)\ne0$. The identity
\[
 v\left(\frac{x}{1+x}\right)=\max(0,v(x))\quad(v(x)\ne0)
\]
shows that every $\Phi_{R_j}(\eta_k)$ is integral.
Consequently $v_p(G_{j,k})\ge3$, and all scalar second summands
in~\eqref{eq:W} have positive valuation at both places over $p$.
No upper bound on the initial formal depth is used.

\subsection{The pole certificates}

The exact contact factorization is
\[
 R_j=\left(\frac{A_j}{v\delta_j}\right)^{10}
                  \prod_{\theta\in\Theta_j}
                         \left(u-(\theta/A_j)v\right).
\]
Thus $R_j$ is a positive rational square times its ten selected
primes, each of which has valuation exactly one and avoids the
support of $\rho$. At the designated place over such a prime,
$W_j^{-1}\equiv1$. At the conjugate place every factor has
nonpositive valuation, since a nonzero second summand has odd
valuation. Its inverse is congruent to zero or one, and so
$W_j^{-1}\equiv0$ or $1$ there. The trace is therefore congruent
to $1$ or $2$ and is a unit. The odd valuation of $b_j$ is positive.
At primes in the support of $\rho$, both $R_j$ and $W_j$ are
arbitrarily close to $1$ by preparation, so the odd valuations
remain nonnegative in $b_j$. At every other prime the valuation
of $\rho R_j$ is even, as is its change under division by the
square trace. Each $b_j$ has square reduced denominator.

At $2,3,17$, choose the same neighborhoods so that $R_j$ is a
local square and $\Tr(W_j^{-1})$ is sufficiently close to $2$.
The squareclass of $b_j$ is then that of $\rho$. At two, the
factor $4$ in~\eqref{eq:poleinputs} cancels the limiting trace
square and gives $v_2(b_j)=4$; the required depths at three and
seventeen are unchanged. At the real embeddings arrange
$1/2<R_j<2$ and $1\le W_j<2$. Every second summand is positive
or zero, and preparation makes their finite product less than two.
The trace is at least one, giving
\[
                              0<b_j<8\rho<17.
\]
All scalar denominators are nonzero by $|R_j\eta_k|<1/2$;
all products and traces are nonzero by positivity. The inherited
converse now supplies the three nonzero pole witnesses in the
specified conic charts. All nine point charts are available, since
the root labels are nonzero and the guarded differences have
$h>2$, $a>1$. This completes the source witness for every integer
input and every fixed odd $L\ge3$.

\section{One polynomial and its numerical bounds}

The retained coordinates consist of the input $X$, eighteen point
coordinates, one nonzero slope, and six pole parameters. Their
number is twenty-six. The finite algebra has nine quartic generators
and three nonzero quadratic pole generators, hence rank
\[
                             4^9 2^3=2^{21}.
\]
The active finite-algebra construction~\cite{Active} eliminates all
twelve finite roots and retains all guards with one additional
rational coordinate. This gives exactly twenty-seven total variables.

Here is its finite polynomial specification. Let $B_0$ be the
algebra of the nine point roots and $A_0$ the algebra after adjoining
the three pole roots $q_1,q_2,q_3$, with $q_j^2=c_j\in B_0$
after rational square clearing. Put $b_0=2^{18}$, $n_0=2^{21}$.
The complete guard has the form $G=gq_1q_2q_3$, with $g\in B_0$.
There are no additional closed equations. Define
\[
 T_0=g^2c_1c_2c_3,\qquad
 \det_{B_0}(I-vM_{T_0})=\sum_{j=0}^{b_0}a_jv^j,\qquad a_0=1.
\]
Let $\theta_1,\ldots,\theta_m$, with $m=12$, be the point and
pole generators, and let $\delta_i$ be the discriminant of point
quartic $i$. All nine discriminants are nonzero at every rational
base specialization. Set
\[
 R=n_0\left(1+\sum_{i=1}^{m}\sum_{u,v=1}^{n_0}
       (M_{\theta_i}^{A_0})_{uv}^2+
       \sum_{u,v=1}^{b_0}(M_g^{B_0})_{uv}^2\right).
\]
For $1\le r\le b_0$, define
\begin{align*}
 c_r&=a_r,&h_r(Z)&=\sum_{j=0}^r a_jZ^{r-j},&E_r&=c_r-h_r(T_0),\\
 D_r&=c_r^2\prod_{i=1}^9\delta_i^2,&s_r&=4r+108,\\
 K_r&=(2R)^{s_r}(1+c_r^2)\prod_{i=1}^9(1+\delta_i^2),&
 C_r&=D_r+4RK_r,\\
 \beta_r&=D_rE_r\left(C_r^m+
             \sum_{i=1}^m D_r^iC_r^{m-i}\theta_i\right),&
 H_r(U)&=\det_{A_0}(I-UM_{\beta_r}).
\end{align*}
Then
\begin{equation}\label{eq:polynomial}
 F_L=H_{b_0}(U)\prod_{r=1}^{b_0-1}
                 \left(H_r(U)^2+\sum_{j=r+1}^{b_0}a_j^2\right)
\end{equation}
is the required polynomial after clearing the fixed integer coefficient
denominators in the prescribed source construction. The active
finite-algebra theorem proves that its rational projection is exactly
the guarded source projection. The displayed product has $262144$
factors. It is a finite specification by integer polynomial operations;
no expanded coefficient table of this very large polynomial is needed.

The independent weighted-degree calculation gives, uniformly in $L$,
\begin{equation}\label{eq:degree}
 \deg F_L\le99950805468739574201057280<2^{87}.
\end{equation}
At the explicit defining member $L=125001$, the maximum absolute
integer coefficient is less than
\begin{equation}\label{eq:coefficients}
                                  2^{\,2^{102}}.
\end{equation}
The separating variable has degree at most $E=1099509530624$.
Writing $D=99950805468739574201057280$, the support is therefore
bounded above by
\[
          \binom{D+27}{27}-\binom{D-E+26}{27}<2^{2198}.
\]
This is a support upper bound, not an expanded monomial count. These bounds, together with
the effective soundness and completeness proofs, establish
Theorem~\ref{thm:main}.

\section{Exact certificates}\label{sec:certificate}

All verification programs use the Python standard library. The
independent constant checker is
\begin{center}
\nolinkurl{computation/few_variables/effective27_v19/}\par
\nolinkurl{check_effective27_constants.py}.
\end{center}
It factors the numerators and denominators of all thirty contacts,
the three poles and all pairwise differences. It verifies the
five extra exceptional primes, their product
\[
                         E_*=1159294788412841845761453254053<2^{100},
\]
the primality and admissibility of $\ell$, all $76249$ input
residues modulo $\ell$, and the strict numerical inequalities
leading to~\eqref{eq:binaryforced}, \eqref{eq:inputlog} and the
contradiction for $\Delta\ne0$. It uses only integers and exact
rational arithmetic, including rational lower bounds for the
exponential where a logarithmic comparison is needed. Its output is
\nolinkurl{effective27-constants-certificate.json}.

The separate source and finite-algebra certificates in the same
module reconstruct the thirty contacts, three scales, sixteen raw
comparison tests, six scalar components, three pole charts and every
guard. The source has $6892$ rational arithmetic nodes, $2587$
integer coefficient nodes and $164$ atomic guard factors. The
source-to-integer conversion, quartic relations, square clearings
and guard coverage are checked by exact symbolic identities.
Fifteen independent reconstructions over three finite fields provide
additional finite evidence for the rational expression formulas;
they are not a proof of universal polynomial identities. The reports
\nolinkurl{independent-source-verification.json} and
\nolinkurl{independent-bound-verification.json} verify the stated
source census, all $262144$ degree cases, the degree and support
bounds, and the coefficient bound at $L=125001$. These finite
checks certify the specification and inequalities stated here; the
arithmetic inputs remain precisely those in Section~1.

\begin{thebibliography}{99}
\interlinepenalty=10000
\bibitem{Active}
\emph{Separating the Nonzero Part of a Finite Algebra},
companion paper, 9 November 2026,
\nolinkurl{active_finite_algebra_20261109_v15.tex}.
\bibitem{Dependencies}
\emph{Arithmetic dependencies of the construction with 29 variables},
companion paper, 8 November 2026,
\nolinkurl{arithmetic_dependencies_20261108_v8.tex}.
\bibitem{Multiplicity}
\emph{Multiplicity of modular factors and quadratic contact-height bounds},
companion paper, 10 November 2026,
\nolinkurl{multiplicity_and_contact_heights_20261110_v19.tex}.
\bibitem{Normalizer}
\emph{Scalar contact normalization by finite-field avoidance},
companion paper, 8 November 2026,
\nolinkurl{scalar_normalization_20261108_v9.tex}.
\bibitem{PointSource}
\emph{A defining polynomial in twenty-seven variables:
an existence theorem with a uniform degree bound},
companion source paper, 9 November 2026,
\nolinkurl{few_variables27_nontorsion_20261109_v15.tex}.
\bibitem{QuadraticCharts}
\emph{A single quartic chart for square roots in a quadratic extension},
companion paper, 8 November 2026,
\nolinkurl{quadratic_square_chart_20261108_v9.tex}.
\bibitem{RawSource}
\emph{A Defining Polynomial in Twenty-Five Variables},
companion paper, 10 November 2026,
\nolinkurl{few_variables25_merged_20261110_v19.tex}.
\bibitem{Source}
E.~Glazer, \emph{Diophantine definition of the integers over the
rationals}, Version~9, \nolinkurl{manuscripts/quartic_v9.tex},
commit \nolinkurl{14e8f601fe0d9173d4a5f084aedc3e07066a7688}.
\url{https://github.com/elliotglazer/diophantine-z-q}.
The inherited restricted-converse correction is identified
in~\cite{Dependencies}; the scalar realization is
\nolinkurl{v5:scalar32-bridge}.
\bibitem{SquareCharts}
\emph{Square abscissas and direct charts for the elliptic subgroup},
companion paper, 8 November 2026,
\nolinkurl{square_abscissa_charts_20261108_v9.tex}.
\bibitem{SquarePredicate}
\emph{Optimal polynomial square tests in quadratic number fields},
companion paper, 8 November 2026,
\nolinkurl{quadratic_square_predicates_20261108_v10.tex}.
\bibitem{Sums}
\emph{Effective heights of sums of elliptic abscissas},
companion paper, 8 November 2026,
\nolinkurl{effective_elliptic_sums_20261108_v10.tex}.
\bibitem{Sun}
Z.-W.~Sun, simultaneous prime-pattern theorem,
Theorem~4.1 of arXiv:2610.20691v2, as pinned in~\cite{Dependencies}.
\end{thebibliography}
\end{document}
